\chapter{Inspection}
\label{ch:u-inspection}

At every stop point of a patrol the robot inspects the \textbf{inspection points} taught for
that stop: it points the thermal camera at each piece of equipment, measures it and judges it,
and listens with the acoustic camera where partial discharge is checked.

\section{Teach the inspection points in NETRA}

\menupath{Substation › Inspection points › New point}:

\begin{enumerate}
  \item Pick the route and its \ui{Stop}, and the type: \ui{Thermal} or \ui{Partial discharge}.
  \item Set the camera's pose: \ui{Pan °}, tilt and zoom, as the camera saw the equipment at that
    stop.
  \item For a thermal point, add a \ui{Target} for each component (phase R, S, T, a clamp) and draw
    its region on the image; give \ui{Max °C}, its rated temperature: a reading at or above it is
    critical (P1).
  \item \ui{Create point}.
\end{enumerate}

The robot takes the plan with every patrol started from NETRA. A point \ui{Not on a route} is
\ui{Taken on request}: NETRA asks the robot to inspect it where it stands (not during a patrol).

\begin{notebox}
The technician can also teach targets on the robot itself, from the camera's current view
(Tuning Guide, ``Inspection''). Both work side by side; NETRA's points are uploaded to NETRA.
\end{notebox}

\section{What happens at a stop}

\begin{center}
\includegraphics[width=\linewidth]{inspection}
\end{center}

A stop with no inspection point just waits its time. A camera that does not answer is
reported, never a reason to end the patrol.

\section{How a finding is judged}

Every thermal target is judged three ways; the finding takes the highest severity of the three:

\begin{center}
\includegraphics[width=\linewidth]{severity}
\end{center}

\begin{center}
\begin{tabularx}{\linewidth}{@{}>{\bfseries}p{0.24\linewidth}X@{}}
  \toprule
  Criterion & Compares \\
  \midrule
  absolute & the hottest region with fixed limits, or the target's rated temperature \\
  ambient & the hottest region with the air temperature of the day \\
  similar & the hottest region with the others of the same target (the three phases) \\
  \bottomrule
\end{tabularx}
\end{center}

\begin{warnbox}[Set the air temperature]
The ambient criterion needs the day's air temperature. The robot uses 30~°C until the
technician sets it; on a hot day, loaded equipment otherwise reads high.
\end{warnbox}

Partial discharge is classified by the acoustic camera:

\begin{center}
\begin{tabularx}{\linewidth}{@{}>{\bfseries}p{0.22\linewidth}p{0.22\linewidth}X@{}}
  \toprule
  Type & Severity & Meaning \\
  \midrule
  noise & normal & no discharge \\
  corona & attention & common outdoors on high-voltage parts: watch the trend \\
  surface & warning & tracking on insulators or terminations \\
  floating & critical & a loose, floating part \\
  \bottomrule
\end{tabularx}
\end{center}

\section{Where the results are}

\begin{itemize}
  \item \textbf{NETRA}: \menupath{Substation › Overview} (every piece of equipment and its last
    capture), \menupath{Findings} (what needs action), each point with its trend of readings and
    its images (thermal, visible, sound map).
  \item \textbf{The H12}: the last point and its severity after \ui{Inspect} or a stop.
  \item \textbf{The robot}: every report with its images and raw thermal frames, by day
    (\file{~/x30/reports}); the oldest days go when the disk runs short.
\end{itemize}

\ui{Inspect} on the H12 inspects the point where the robot stands, now: useful to check a
newly taught point. It takes a few minutes for a point with several targets.
